\section{Math}
\subsection{Basis Vector}
\begin{minted}{c++}
template<typename T = int, int LG = 31>
struct Basis {
  T basis[LG]; int sz;
  vector<int> who;
  Basis() { reset(); }
  void reset() {
    fill(basis, basis + LG, 0);
    sz = 0; who.clear();
  }
  bool insert(T mask, int idx = -1) {
    for (int i = LG - 1; i >= 0; --i) {
      if (mask >> i & 1) {
        if (basis[i]) mask ^= basis[i];
        else { basis[i] = mask; ++sz;
          who.push_back(idx); return true; }
      }
    }
    return false;
  }
  T maxXor() {
    T ans = 0;
    for (int i = LG - 1; i >= 0; --i)
      ans = max(ans, ans ^ basis[i]);
    return ans;
  }
  bool can(T x) {
    for (int i = LG - 1; i >= 0; --i) {
      if (x >> i & 1) {
        if (!basis[i]) return false;
        x ^= basis[i];
      }
    }
    return true;
  }
  T kth(T k) {
    if (k < 1 || k > ((T)1 << sz)) return -1;
    T mask = 0;
    int64_t tot = 1LL << sz;
    for (int i = LG - 1; i >= 0; --i) {
      if (basis[i]) {
        int64_t low = tot / 2;
        if ((low < k && (~mask >> i & 1)) ||
            (low >= k && (mask >> i & 1)))
          mask ^= basis[i];
        if (low < k) k -= low;
        tot /= 2;
      }
    }
    return mask;
  }
  T count_less(T x) {
    if (x < 0) return 0;
    T ret = 0, cnt = ((T)1 << sz), mask = 0;
    for (int i = LG - 1; i >= 0; i--) {
      if (basis[i]) {
        if (x >> i & 1) {
          ret += (cnt >> 1);
          if (~mask >> i & 1) mask ^= basis[i];
        } else {
          if (mask >> i & 1) mask ^= basis[i];
        }
        cnt >>= 1;
      } else {
        if ((x ^ mask) >> i & 1) {
          if (x >> i & 1) return ret + cnt;
          else return ret;
        }
      }
    }
    return ret;
  }
  T count_leq(T x) { return count_less(x + 1); }
  T jumpK(T val, T k) {
    return kth(count_leq(val) + k);
  }
};
\end{minted}
\subsection{Matrix Exponentiation}
\begin{minted}{c++}
vector<vector<int>> multiply(
  vector<vector<int>>& a,
  vector<vector<int>>& b) {
  int n = sz(a), m = sz(a[0]), p = sz(b[0]);
  vector<vector<int>> result(n,
    vector<int>(p, 0));
  for (int i = 0; i < n; i++)
    for (int k = 0; k < m; k++) {
      if (a[i][k] == 0) continue;
      for (int j = 0; j < p; j++)
        result[i][j] = ((result[i][j] +
          a[i][k] * b[k][j]) % mod + mod) % mod;
    }
  return result;
}
vector<vector<int>> expo(
  vector<vector<int>> t, int p) {
  int n = t.size();
  vector<vector<int>> result(n, vector<int>(n, 0));
  for (int i = 0; i < n; i++) result[i][i] = 1;
  while (p > 0) {
    if (p & 1) result = multiply(result, t);
    t = multiply(t, t);
    p >>= 1;
  }
  return result;
}
\end{minted}
\subsection{Linear Combination Checking}
Given array \texttt{arr} and integer N, check if N is
a linear combination of arr's elements:
$N = a \cdot x + b \cdot y + c \cdot z + \ldots$

Ans: $N \bmod \text{GCD}(\text{all elements}) == 0$.

\subsection{Math Formula Cheatsheet}
\textbf{Geometric Series}: \\
\textbullet{} $\sum_{i=a}^b c^i = \frac{c^{b+1} - c^a}{c - 1}$ $(c \neq 1)$ \\
\textbullet{} $\sum_{i=0}^n c^i = \frac{c^{n+1}-1}{c-1}$ $(c \neq 1)$ \\
\textbullet{} $\sum_{i=0}^\infty c^i = \frac{1}{1-c}$ $(|c| < 1)$

\textbf{Sums of Powers}: \\
\textbullet{} $\sum_{i=1}^n i = \frac{n(n+1)}{2}$ \\
\textbullet{} $\sum_{i=1}^n i^2 = \frac{n(n+1)(2n+1)}{6}$ \\
\textbullet{} $\sum_{i=1}^n i^3 = \frac{n^2(n+1)^2}{4}$ \\
\textbullet{} $\sum_{i=1}^n i^4 = \frac{n(n+1)(2n+1)(3n^2+3n-1)}{30}$

\textbf{Harmonic Sums}: \\
\textbullet{} $H_n = \sum_{i=1}^n \frac{1}{i}$ \\
\textbullet{} $\sum_{i=1}^n i H_i = \frac{n(n+1)}{2} H_n - \frac{n(n-1)}{4}$

\textbf{Combinatorial Sum}: \\
\textbullet{} $\sum_{k=0}^n \binom{r+k}{k} = \binom{r+n+1}{n}$
